\chapter{Quantum computation: gates, algorithms, and protocols}
\label{ch:quantum-information-processing}

Parts~I--IV built a physical qubit from the ground up: a two-level
system from the postulates (\cref{ch:postulates,ch:model-systems}), its
decoherence from system--bath coupling (\cref{ch:open-systems}), and a
concrete superconducting realisation with gates and readout
(\cref{ch:circuit-qed}). This chapter steps up one level of
abstraction. We forget, for the moment, what the qubit is made of and
ask what can be \emph{done} with it as a carrier of information:
how unitary gates compose into circuits, what set of gates suffices to
build any computation, and what tasks become possible that have no
classical analogue. Three threads run through the chapter and recur
later. First, computation is reversible unitary evolution punctuated by
measurement --- the circuit model. Second, the impossibility of copying
an unknown state (the no-cloning theorem) is not a limitation but the
structural fact underwriting teleportation and quantum cryptography.
Third, not all quantum gates are equal: the Clifford gates are
classically simulable, and the computational ``hardness'' --- the part
that resists classical simulation and that error correction must pay
for --- lives in a thin layer of non-Clifford resources. That last
observation is the bridge to quantum error correction
(\cref{ch:qec}), and the stabiliser language we develop here is exactly
the language that chapter speaks.

% --------------------------------------------------------------
\section{The circuit model and a universal gate set}
\label{sec:circuit-model}

\begin{phenomenon}
A classical computer reduces any computation to a handful of logic
gates --- \textsc{and}, \textsc{or}, \textsc{not} --- wired together.
Most of these gates are irreversible: \textsc{and} maps two input bits
to one output and erases information. A quantum computer is built the
same way, from a small set of elementary gates wired into a circuit,
but with one structural difference forced by physics. Closed-system
evolution is unitary (\cref{ch:postulates}), and unitaries are
invertible and norm-preserving. Quantum logic is therefore
\emph{reversible}: every gate has an inverse, and information is never
destroyed until a measurement is made. Output is read out only at the
end, by measurement, which collapses the state and returns bits.
\end{phenomenon}

\begin{definition}[The quantum circuit model]
\label{def:circuit-model}
A register of $n$ qubits is a state in
$(\C^2)^{\otimes n}$. A \emph{quantum circuit} is a finite, ordered
product of \emph{gates} --- unitaries each acting on one or two qubits
and extended by the identity on the rest --- applied to an initial
product state (conventionally $\ket{0}^{\otimes n}$), followed by
measurement of some qubits in the computational basis. In the standard
diagrammatic convention each qubit is a horizontal wire, time flows
left to right, a boxed symbol $\op U$ on a wire is a single-qubit gate,
a filled control dot joined by a vertical line to a target
($\oplus$ for a bit flip) is a controlled gate, and a meter symbol
denotes a computational-basis measurement. We write circuits
algebraically as the corresponding ordered operator product, read
right to left.
\end{definition}

The elementary single-qubit gates are the Pauli operators of
\cref{ch:model-systems} --- the bit flip $\hat\sigma_x$, the phase flip
$\hat\sigma_z$, and $\hat\sigma_y=i\hat\sigma_x\hat\sigma_z$ --- together
with three gates that have no classical counterpart: the
\emph{Hadamard} $\mathrm{H}$, the \emph{phase gate} $\mathrm{S}$, and
the \emph{$\mathrm{T}$ gate},
\begin{equation}
\mathrm{H}=\frac{1}{\sqrt2}\begin{pmatrix}1&1\\1&-1\end{pmatrix},
\qquad
\mathrm{S}=\begin{pmatrix}1&0\\0&i\end{pmatrix},
\qquad
\mathrm{T}=\begin{pmatrix}1&0\\0&e^{i\pi/4}\end{pmatrix},
\end{equation}
together with the continuous rotations generated by the Paulis,
\begin{equation}
\label{eq:rotations}
\op R_{\hat n}(\theta)
=\exp\!\Bigl(-i\tfrac{\theta}{2}\,\hat n\cdot\hat{\vec\sigma}\Bigr)
=\cos\tfrac{\theta}{2}\,\id
-i\sin\tfrac{\theta}{2}\,(\hat n\cdot\hat{\vec\sigma}),
\end{equation}
which are the abstract form of the hardware single-qubit gates of
\cref{stmt:single-qubit-gates}. The Hadamard builds superpositions,
$\mathrm{H}\ket{0}=\ket{+}$ and $\mathrm{H}\ket{1}=\ket{-}$, and is its
own inverse, $\mathrm{H}^2=\id$. Note $\mathrm{S}=\mathrm{T}^2$ and
$\mathrm{S}^2=\hat\sigma_z$. The canonical two-qubit gate is the
\emph{controlled-NOT} of \cref{ch:composite},
\begin{equation}
\mathrm{CNOT}
=\ket{0}\!\bra{0}\otimes\id+\ket{1}\!\bra{1}\otimes\hat\sigma_x
=\begin{pmatrix}1&0&0&0\\0&1&0&0\\0&0&0&1\\0&0&1&0\end{pmatrix},
\qquad
\mathrm{CZ}=\mathrm{diag}(1,1,1,-1),
\end{equation}
with $\mathrm{CZ}=(\id\otimes\mathrm{H})\,\mathrm{CNOT}\,(\id\otimes\mathrm{H})$
relating the two by a Hadamard on the target.

\begin{statement}[\textbf{Universality of single-qubit gates with CNOT}~\cite{nielsen2010}]
\label{stmt:universal-gates}
Any unitary on $n$ qubits can be written exactly as a finite circuit of
single-qubit gates and CNOTs. Moreover the discrete set
$\{\mathrm{H},\mathrm{T},\mathrm{CNOT}\}$ is \emph{universal}: it
generates a dense subgroup of the $n$-qubit unitaries, so any target
unitary can be approximated to operator-norm error $\varepsilon$ by a
finite circuit. The Solovay--Kitaev theorem guarantees this costs only
$O\!\bigl(\log^{c}(1/\varepsilon)\bigr)$ gates per single-qubit
unitary, with $c\approx 2$.
\end{statement}

\begin{proof}[Derivation (sketch)]
The argument has three layers. \emph{(i) Exact decomposition into
two-level rotations.} Any $d\times d$ unitary is a product of unitaries
each acting nontrivially on a two-dimensional subspace (a Givens
rotation); for $d=2^n$ there are $O(d^2)$ of them. \emph{(ii) Two-level
unitaries from controlled gates.} Each two-level rotation, after a Gray-code
relabelling of basis states, is a single-qubit rotation controlled on
the remaining $n-1$ qubits; a controlled-$\op U$ with many controls
reduces to a sequence of CNOTs and singly-controlled gates. The
singly-controlled $\op U$ itself factorises: writing
$\op U=e^{i\alpha}\op A\,\hat\sigma_x\,\op B\,\hat\sigma_x\,\op C$ with
$\op A\op B\op C=\id$ (always possible by the Euler decomposition of
$SU(2)$), the controlled-$\op U$ is implemented by a phase $e^{i\alpha}$
on the control followed by $\op C$, CNOT, $\op B$, CNOT, $\op A$ on the
target --- single-qubit gates and two CNOTs. \emph{(iii)
Discretisation.} The remaining single-qubit rotations are approximated
from $\{\mathrm{H},\mathrm{T}\}$: $\mathrm{T}$ is a rotation by the
irrational-of-$\pi$ angle $\pi/4$ about $\hat z$, and
$\mathrm{H}\mathrm{T}\mathrm{H}$ rotates about $\hat x$; two rotations
about non-parallel axes by an angle incommensurate with $2\pi$ generate
a dense subgroup of $SU(2)$. Solovay--Kitaev makes the approximation
efficient.
\end{proof}

The Clifford gates $\{\mathrm{H},\mathrm{S},\mathrm{CNOT}\}$ generate an
important \emph{finite} subgroup, the Clifford group, which we return to
in \cref{sec:magic-states}; adding the single non-Clifford gate
$\mathrm{T}$ tips the set over into universality.

\begin{example}[Making a Bell pair: the Hadamard--CNOT circuit]
\label{ex:bell-circuit}
The smallest nontrivial circuit prepares the Bell state of
\cref{ex:bell-states}. Start in $\ket{00}$, apply a Hadamard to the
first qubit, then a CNOT with the first as control:
\begin{align}
\ket{00}
&\xrightarrow{\;\mathrm{H}\otimes\id\;}
\ket{+}\ket{0}
=\tfrac{1}{\sqrt2}\bigl(\ket{0}+\ket{1}\bigr)\ket{0}\\
&\xrightarrow{\;\mathrm{CNOT}\;}
\tfrac{1}{\sqrt2}\bigl(\ket{00}+\ket{11}\bigr)=\ket{\Phi^+}.
\end{align}
The Hadamard creates the superposition, and the CNOT correlates it ---
generating exactly one ebit (\cref{ch:composite}) from a product input.
This two-gate primitive is the entangling front end of nearly every
protocol below, and its hardware incarnations (cross-resonance, iSWAP-,
and controlled-phase gates) are the subject of
\cref{sec:two-qubit-gates}.
\end{example}

\paragraph{Transition.}
The circuit model gives us a finite vocabulary --- a handful of gates
--- with which to build any unitary. Before exploiting it, we pause on
a property of quantum information that has no classical analogue and
that turns out to be the hinge of the rest of the chapter: unknown
states cannot be copied.

% --------------------------------------------------------------
\section{The no-cloning theorem}
\label{sec:no-cloning}

\begin{phenomenon}
Classical information can be copied at will: a file is duplicated bit by
bit, and the copy is indistinguishable from the original. Could a
machine do the same for an unknown qubit --- take in one copy of
$\ket{\psi}$ and put out two? If it could, quantum mechanics would
unravel in several places at once. One could defeat the measurement
problem by cloning a state many times and measuring each copy to read
off $\ket{\psi}$ exactly; one could signal faster than light by cloning
half of an entangled pair. No such device has ever been built, and the
reason is a one-line theorem.
\end{phenomenon}

\begin{theorem}[No-cloning~\cite{wootters1982}]
\label{thm:no-cloning}
There is no unitary $\op U$ on $\hilbert\otimes\hilbert$ and fixed
``blank'' state $\ket{s}$ such that
$\op U\bigl(\ket{\psi}\otimes\ket{s}\bigr)=\ket{\psi}\otimes\ket{\psi}$
for every $\ket{\psi}\in\hilbert$.
\end{theorem}

\begin{proof}[Derivation]
Suppose such a $\op U$ existed and acted as required on two states
$\ket{\psi}$ and $\ket{\phi}$,
\begin{equation}
\op U\bigl(\ket{\psi}\ket{s}\bigr)=\ket{\psi}\ket{\psi},
\qquad
\op U\bigl(\ket{\phi}\ket{s}\bigr)=\ket{\phi}\ket{\phi}.
\end{equation}
Take the inner product of the two left-hand sides and of the two
right-hand sides. Unitarity preserves inner products, so the left-hand
sides give
$\bra{\phi}\bra{s}\op U^\dagger\op U\ket{\psi}\ket{s}
=\braket{\phi}{\psi}\braket{s}{s}=\braket{\phi}{\psi}$,
while the right-hand sides give
$\braket{\phi}{\psi}\braket{\phi}{\psi}=\braket{\phi}{\psi}^2$.
Hence
\begin{equation}
\braket{\phi}{\psi}=\braket{\phi}{\psi}^2
\quad\Longrightarrow\quad
\braket{\phi}{\psi}\in\{0,1\}.
\end{equation}
The overlap can only be $0$ (orthogonal) or $1$ (identical). A cloner
therefore cannot exist for a general pair of non-orthogonal, distinct
states. The same conclusion follows from linearity: a device that
clones $\ket{0}$ and $\ket{1}$ acts on $\ket{+}\ket{s}$ to give
$\tfrac{1}{\sqrt2}(\ket{00}+\ket{11})$, not the required
$\ket{+}\ket{+}$.
\end{proof}

\begin{remark}
The theorem forbids copying \emph{unknown superpositions}, not copying
basis states. A CNOT does copy the computational basis,
$\mathrm{CNOT}\ket{a}\ket{0}=\ket{a}\ket{a}$ for $a\in\{0,1\}$ --- this
is ordinary classical copying of a known bit. What fails is copying a
state in superposition: the same CNOT entangles rather than duplicates
(\cref{thm:no-cloning} derivation, last line).
\end{remark}

\begin{application}[Why no-cloning matters]
\label{app:no-cloning-uses}
Three consequences thread through the rest of this part. (i) An
eavesdropper cannot copy qubits in transit to measure them at leisure;
this is the structural basis of quantum key distribution
(\cref{sec:qkd}). (ii) Teleportation (\cref{sec:teleportation}) moves a
state without violating no-cloning precisely because the original is
destroyed by Alice's measurement --- there is never a second copy.
(iii) Entanglement cannot be used for faster-than-light signalling: if
it could, cloning would let one read the signal, so the no-signalling
principle and no-cloning stand or fall together. The theorem is also the
reason quantum error correction (\cref{ch:qec}) cannot work by the
classical trick of keeping redundant copies, forcing the subtler
redundancy of \cref{sec:magic-states}'s stabiliser codes.
\end{application}

\paragraph{Transition.}
No-cloning sounds purely restrictive, but it is the enabling condition
for the two dual communication protocols we treat next. Both spend the
entanglement resource quantified in \cref{ch:composite} --- one ebit
per use --- and together they pin down the operational meaning of that
unit.

% --------------------------------------------------------------
\section{Teleportation and superdense coding}
\label{sec:teleportation}

\begin{phenomenon}
Alice holds a qubit in an unknown state $\ket{\psi}$ and wants Bob to
have it, but they share only a classical telephone line and, in
advance, one entangled pair. Can she send the qubit without any quantum
channel? The dual question: Alice and Bob again share one entangled
pair, and Alice may send Bob a single physical qubit; how many
\emph{classical} bits can she thereby convey? The answers --- yes, at a
cost of two classical bits (teleportation); and two classical bits per
qubit (superdense coding) --- are mirror images, and they make the ebit
of \cref{ch:composite} a concrete, expendable currency.
\end{phenomenon}

\begin{statement}[\textbf{Quantum teleportation}~\cite{bennett1993}]
\label{stmt:teleportation}
One ebit of pre-shared entanglement plus two classical bits suffice to
transmit one unknown qubit. No physical qubit travels from sender to
receiver, and the sender's copy is destroyed in the process, consistent
with \cref{thm:no-cloning}.
\end{statement}

\begin{proof}[Derivation]
Label Alice's unknown qubit $C$ and the shared pair $A$ (Alice) and $B$
(Bob), prepared in $\ket{\Phi^+}_{AB}$. The full state, with
$\ket{\psi}=\alpha\ket{0}+\beta\ket{1}$, is
\begin{equation}
\ket{\psi}_C\otimes\ket{\Phi^+}_{AB}
=\tfrac{1}{\sqrt2}\bigl[
\alpha\ket{0}_C(\ket{00}+\ket{11})_{AB}
+\beta\ket{1}_C(\ket{00}+\ket{11})_{AB}\bigr].
\end{equation}
Alice applies $\mathrm{CNOT}_{C\to A}$, which flips $A$ when $C=1$, then
a Hadamard on $C$, which splits the $C$ kets:
\begin{align}
\xrightarrow{\;\mathrm{CNOT}_{C\to A}\;}\ \tfrac{1}{\sqrt2}\bigl[
&\alpha\ket{0}_C(\ket{00}+\ket{11})_{AB}
+\beta\ket{1}_C(\ket{10}+\ket{01})_{AB}\bigr],\\
\xrightarrow{\;\mathrm{H}_C\;}\ \tfrac12\bigl[
&\alpha(\ket{0}+\ket{1})_C(\ket{00}+\ket{11})_{AB}\notag\\
&{}+\beta(\ket{0}-\ket{1})_C(\ket{10}+\ket{01})_{AB}\bigr].
\end{align}
Regrouping by the four possible outcomes $(m_C,m_A)$ of measuring $C$
and $A$ in the computational basis,
\begin{equation}
\begin{aligned}
={}\tfrac12\bigl[
&\ket{00}_{CA}(\alpha\ket{0}+\beta\ket{1})_B
+\ket{01}_{CA}(\alpha\ket{1}+\beta\ket{0})_B\\
&{}+\ket{10}_{CA}(\alpha\ket{0}-\beta\ket{1})_B
+\ket{11}_{CA}(\alpha\ket{1}-\beta\ket{0})_B\bigr].
\end{aligned}
\end{equation}
Each outcome leaves Bob's qubit in
$\hat\sigma_x^{\,m_A}\hat\sigma_z^{\,m_C}\ket{\psi}$. Alice sends the two
bits $(m_C,m_A)$ over the classical channel, and Bob undoes the Pauli by
applying $\hat\sigma_z^{\,m_C}\hat\sigma_x^{\,m_A}$, recovering
$\ket{\psi}$ exactly. Two facts close the loop: Alice's qubit $C$ is now
a measured classical bit, so no clone exists (\cref{thm:no-cloning});
and before Bob learns $(m_C,m_A)$ his reduced state is the maximally
mixed $\tfrac12\id$, identical for every input, so no information has
travelled faster than the classical bits --- there is no
faster-than-light signalling.
\end{proof}

\begin{statement}[\textbf{Superdense coding}~\cite{bennett1992}]
\label{stmt:superdense}
One ebit of pre-shared entanglement lets one transmitted qubit carry
two classical bits. It is the time-reverse of teleportation: the same
ebit, with the roles of quantum and classical channels exchanged.
\end{statement}

\begin{proof}[Derivation]
Alice and Bob share $\ket{\Phi^+}_{AB}$. To send the two bits
$(b_1,b_2)$, Alice applies $\hat\sigma_z^{\,b_1}\hat\sigma_x^{\,b_2}$ to
her half $A$, mapping the shared state onto one of the four Bell states,
\begin{equation}
\begin{aligned}
(0,0)&\to\ket{\Phi^+}, &\qquad
(0,1)&\to\hat\sigma_x^{(A)}\ket{\Phi^+}=\ket{\Psi^+},\\
(1,0)&\to\hat\sigma_z^{(A)}\ket{\Phi^+}=\ket{\Phi^-}, &\qquad
(1,1)&\to\hat\sigma_z^{(A)}\hat\sigma_x^{(A)}\ket{\Phi^+}=\ket{\Psi^-}.
\end{aligned}
\end{equation}
She then sends qubit $A$ to Bob, who now holds both halves. Because the
four Bell states are orthonormal, Bob distinguishes them perfectly with
a Bell measurement --- $\mathrm{CNOT}_{A\to B}$, then $\mathrm{H}$ on
$A$, then measure both --- reading off $(b_1,b_2)$ without error. A
single qubit thus delivers two classical bits, but only with the ebit
spent in advance: a lone qubit, with no shared entanglement, can carry
at most one bit of accessible information.
\end{proof}

\begin{application}[The ebit as currency]
\label{app:ebit-currency}
The two protocols give the ebit of \cref{ch:composite} an exchange
rate. Teleportation converts
\(\text{1 ebit}+\text{2 cbits}\to\text{1 qubit}\); superdense coding
converts \(\text{1 ebit}+\text{1 qubit}\to\text{2 cbits}\). Reading them
together shows the resources are not freely interchangeable --- a qubit
is strictly more than a classical bit, and entanglement is the catalyst
that lets one be traded for the other. This accounting is the reason
entanglement is treated as a quantifiable resource throughout quantum
information, and it is the same entanglement that
\cref{ch:coherent-multiqubit} exploits to beat the standard quantum
limit in sensing.
\end{application}

\paragraph{Transition.}
Teleportation relied on the fact that Eve, intercepting the channel,
could learn nothing without the classical bits; no-cloning forbade her
from hoarding copies. Turned around, these same facts let two parties
detect any eavesdropper and so distil a provably secret key --- the
subject of quantum cryptography.

% --------------------------------------------------------------
\section{Quantum key distribution}
\label{sec:qkd}

\begin{phenomenon}
Two parties want to agree on a secret string of bits --- a key for
encryption --- over a channel an eavesdropper may fully control.
Classical key exchange rests on the presumed hardness of a mathematical
problem (factoring, discrete logarithms), which a quantum computer would
break. Quantum key distribution rests instead on physical law: any
attempt to learn the key disturbs the quantum states carrying it, and
the disturbance is detectable. Security comes from no-cloning and the
unavoidable back-action of measurement, not from computational
assumptions.
\end{phenomenon}

\begin{statement}[\textbf{BB84 and E91 key distribution}~\cite{bennett1984,ekert1991}]
\label{stmt:bb84}
In \emph{BB84}, Alice encodes random bits in one of two conjugate bases
chosen at random --- the $\hat\sigma_z$ basis $\{\ket{0},\ket{1}\}$ or
the $\hat\sigma_x$ basis $\{\ket{+},\ket{-}\}$ --- and Bob measures each
qubit in a randomly chosen basis. Keeping only the rounds where their
bases agree (sifting) yields a shared key; comparing a public sample
bounds the eavesdropper's information through the induced error rate. In
the entanglement-based \emph{E91} variant, Alice and Bob share Bell
pairs and certify security directly by testing the CHSH inequality
(\cref{sec:quantum-randomness}): a genuine quantum correlation
saturating the Tsirelson bound leaves no room for a hidden eavesdropper.
\end{statement}

\begin{proof}[Derivation: why eavesdropping shows up]
Consider the simplest attack, intercept-and-resend: Eve measures each
qubit and forwards a fresh qubit in the state she found. Suppose Alice
sends $\ket{+}$ (a $\hat\sigma_x$-basis bit) but Eve, not knowing the
basis, measures in the $\hat\sigma_z$ basis. She obtains $\ket{0}$ or
$\ket{1}$ with equal probability and resends it. When Bob later measures
in the correct $\hat\sigma_x$ basis, the resent
$\hat\sigma_z$-eigenstate gives $\ket{+}$ or $\ket{-}$ with equal
probability,
\begin{equation}
|\!\braket{-}{0}\!|^2=|\!\braket{-}{1}\!|^2=\tfrac12,
\end{equation}
so Bob's bit disagrees with Alice's half the time on these rounds. Eve
chooses the wrong basis on half of all sifted rounds, so the
intercept-resend attack injects a bit-error rate of
$\tfrac12\times\tfrac12=25\%$. No-cloning (\cref{thm:no-cloning}) blocks
the obvious workaround --- copying each qubit to measure later in the
right basis --- because the required cloner does not exist. Alice and
Bob estimate the error rate from a public subset of the key; an
anomalously high rate signals an eavesdropper and the key is discarded.
Below a threshold ($\approx 11\%$ for BB84) the residual leaked
information can be removed by classical error correction and privacy
amplification, leaving a provably secret key. In E91 the same logic is
phrased through the CHSH quantity $S$: undisturbed Bell pairs reach the
quantum value $S=2\sqrt2$ (\cref{sec:quantum-randomness}), whereas any
intercepting strategy that extracts which-path information necessarily
pulls $S$ back towards the classical bound \cref{eq:chsh-bound},
$|S|\le2$, betraying its presence.
\end{proof}

\begin{application}[Security from physics]
\label{app:qkd}
BB84 and E91 turn the foundational results of \cref{ch:postulates}
into a working technology: conjugate observables, measurement
disturbance, and Bell-inequality violation become the guarantees of a
cryptographic protocol now deployed over optical fibre and
satellite links. The shift is conceptual as much as practical --- the
security proof does not assume the eavesdropper is computationally
limited, only that quantum mechanics is correct.
\end{application}

\paragraph{Transition.}
Teleportation, dense coding, and key distribution all used a small fixed
repertoire --- Hadamards, CNOTs, Pauli corrections, basis measurements.
Strikingly, a quantum computer restricted to \emph{exactly} this
repertoire is no more powerful than a classical one. Understanding why
isolates the resource that makes quantum computation hard to simulate,
and points directly at what error correction must protect.

% --------------------------------------------------------------
\section{Clifford gates, the Gottesman--Knill theorem, and magic states}
\label{sec:magic-states}

\begin{phenomenon}
One expects every quantum circuit to be hard to simulate classically ---
that is the premise of quantum speed-up. Yet a large and natural class
of circuits, including the entangling Hadamard--CNOT networks behind
Bell pairs, teleportation, and dense coding, can be tracked on a
classical computer in polynomial time. The quantum advantage cannot
live in those gates. It must live in whatever is left over once they are
removed --- a surprisingly thin layer of ``magic'' that turns out to be
both the source of computational power and the most expensive thing to
protect in a fault-tolerant machine.
\end{phenomenon}

The \emph{Pauli group} on $n$ qubits consists of tensor products of
$\{\id,\hat\sigma_x,\hat\sigma_y,\hat\sigma_z\}$ with an overall phase.
The \emph{Clifford group} is its normaliser: the unitaries $\op C$ that
map Paulis to Paulis under conjugation,
$\op C\,\hat P\,\op C^\dagger\in\text{Pauli group}$. It is generated by
$\{\mathrm{H},\mathrm{S},\mathrm{CNOT}\}$, and the generators act on the
single-qubit Paulis by
\begin{equation}
\label{eq:clifford-conjugation}
\mathrm{H}:\;
\hat\sigma_x\!\leftrightarrow\!\hat\sigma_z,\;
\hat\sigma_y\!\to\!-\hat\sigma_y;
\qquad
\mathrm{S}:\;
\hat\sigma_x\!\to\!\hat\sigma_y,\;
\hat\sigma_y\!\to\!-\hat\sigma_x,\;
\hat\sigma_z\!\to\!\hat\sigma_z;
\end{equation}
while CNOT (control $1$, target $2$) sends
$\hat\sigma_x^{(1)}\to\hat\sigma_x^{(1)}\hat\sigma_x^{(2)}$ and
$\hat\sigma_z^{(2)}\to\hat\sigma_z^{(1)}\hat\sigma_z^{(2)}$, leaving
$\hat\sigma_x^{(2)}$ and $\hat\sigma_z^{(1)}$ fixed.

\begin{statement}[\textbf{Gottesman--Knill theorem}~\cite{gottesman1997}]
\label{stmt:gottesman-knill}
A circuit composed only of Clifford gates, computational-basis state
preparation, and measurement of Pauli observables can be simulated on a
classical computer in time polynomial in the number of qubits and
gates. Clifford gates alone are therefore not universal and yield no
quantum speed-up.
\end{statement}

\begin{proof}[Derivation]
A \emph{stabiliser state} of $n$ qubits is specified not by its $2^n$
amplitudes but by $n$ commuting Pauli operators that fix it (its
\emph{stabilisers}); $\ket{0}^{\otimes n}$ is stabilised by
$\hat\sigma_z^{(1)},\dots,\hat\sigma_z^{(n)}$. Each generator is an
$n$-qubit Pauli, recorded in $O(n)$ bits, so the whole description is
$O(n^2)$ bits. A Clifford gate $\op C$ maps the state to another
stabiliser state whose generators are $\op C\hat P_i\op C^\dagger$ ---
again Paulis, by definition of the Clifford group --- and updating all
$n$ generators costs $O(n^2)$ work per gate using the rules
\cref{eq:clifford-conjugation}. Pauli measurements update the generators
by the same bookkeeping. The entire computation is therefore tracked
classically in polynomial time, never touching the exponentially large
amplitude vector. (This stabiliser bookkeeping is exactly the formalism
on which the codes of \cref{ch:qec} are built.)
\end{proof}

\begin{statement}[\textbf{Magic states restore universality}~\cite{bravyi2005}]
\label{stmt:magic-distillation}
Augmenting the Clifford gates with any non-Clifford gate yields a
universal set. Equivalently, the non-Clifford gate need never be
performed directly: a $\mathrm{T}$ gate can be realised with Clifford
operations and measurement alone, provided one consumes a \emph{magic
state} $\ket{\mathrm{T}}=\mathrm{T}\ket{+}
=\tfrac{1}{\sqrt2}(\ket{0}+e^{i\pi/4}\ket{1})$. Many noisy magic states
can be purified into fewer high-fidelity ones by a Clifford-only
\emph{magic-state distillation} protocol, so the non-Clifford resource
is manufactured offline.
\end{statement}

\begin{proof}[Derivation: the $\mathrm{T}$ gadget]
That $\mathrm{T}$ is non-Clifford is a one-line check:
\begin{equation}
\mathrm{T}\,\hat\sigma_x\,\mathrm{T}^\dagger
=\tfrac{1}{\sqrt2}\bigl(\hat\sigma_x+\hat\sigma_y\bigr),
\end{equation}
which is not a Pauli. To apply it without ever running it directly, take
the data qubit $\ket{\psi}=a\ket{0}+b\ket{1}$ and a magic state
$\ket{\mathrm{T}}$, and apply $\mathrm{CNOT}$ with the data as control
and the magic state as target:
\begin{equation}
\begin{aligned}
(a\ket{0}+b\ket{1})&\otimes\tfrac{1}{\sqrt2}(\ket{0}+e^{i\pi/4}\ket{1})\\
\xrightarrow{\;\mathrm{CNOT}\;}\ \tfrac{1}{\sqrt2}\bigl[
&(a\ket{0}+be^{i\pi/4}\ket{1})\ket{0}
+(ae^{i\pi/4}\ket{0}+b\ket{1})\ket{1}\bigr].
\end{aligned}
\end{equation}
Measuring the second qubit gives outcome $0$, leaving the data in
$a\ket{0}+be^{i\pi/4}\ket{1}=\mathrm{T}\ket{\psi}$; or outcome $1$,
leaving it in $\mathrm{T}^\dagger\ket{\psi}$ up to a global phase, which
the Clifford correction $\mathrm{S}$ converts to $\mathrm{T}\ket{\psi}$.
Every operation performed --- CNOT, measurement, and the $\mathrm{S}$
correction --- is Clifford; all of the non-Clifford character was
imported from the consumed state $\ket{\mathrm{T}}$. This is why the
hard part of a computation can be cleanly separated and outsourced to a
supply of magic states.
\end{proof}

\begin{application}[Where the cost goes in a fault-tolerant computer]
\label{app:magic-cost}
The Clifford/non-Clifford split shapes the architecture of
\cref{ch:qec}. In a good error-correcting code the Clifford gates often
admit a \emph{transversal} (and so naturally fault-tolerant)
implementation, but the Eastin--Knill theorem~\cite{eastin2009}
forbids \emph{any} code from having a universal transversal gate set ---
so at least one gate, conventionally $\mathrm{T}$, cannot be done
transversally. The resolution is exactly \cref{stmt:magic-distillation}:
$\mathrm{T}$ gates are supplied by distilled magic states, and in
realistic resource estimates the magic-state factories dominate the
qubit and time budget. The thin layer of ``magic'' isolated here is thus
also the dominant cost of fault tolerance --- a thread we pick up in the
next chapter.
\end{application}

% --------------------------------------------------------------
This chapter abstracted the qubit into an information carrier. The
circuit model reduced computation to a finite gate vocabulary, and we
saw that single-qubit gates with CNOT --- or the discrete set
$\{\mathrm{H},\mathrm{T},\mathrm{CNOT}\}$ --- can build any unitary. The
no-cloning theorem, far from a mere restriction, underwrote
teleportation, superdense coding, and the physical security of quantum
key distribution. Finally, the Gottesman--Knill theorem split the gate
set in two: the classically simulable Clifford layer, and the thin
non-Clifford layer of magic states that carries the computational power.

That last split is the natural entry to error correction. Everything in
this chapter assumed perfect, noiseless gates and qubits --- but the
qubits of Parts~II and IV decohere, and a real circuit accumulates
errors at every step. The next chapter asks how a logical qubit can be
protected from that noise, and answers it in the very stabiliser
language introduced in \cref{stmt:gottesman-knill}: error correction is
Clifford bookkeeping turned into a shield, with the magic states of
\cref{stmt:magic-distillation} paying for the universality the shield
would otherwise forbid.

\clearpage
\begin{chaptersummary}

\paragraph{Circuit model and universality.}
$n$ qubits live in $(\C^2)^{\otimes n}$; a circuit is an ordered product
of one- and two-qubit unitaries followed by measurement. Single-qubit
gates plus CNOT are exactly universal; the discrete set
$\{\mathrm{H},\mathrm{T},\mathrm{CNOT}\}$ is approximately universal,
with Solovay--Kitaev cost $O(\log^c(1/\varepsilon))$. The
Hadamard--CNOT pair makes a Bell state,
$\mathrm{CNOT}(\mathrm{H}\otimes\id)\ket{00}=\ket{\Phi^+}$.

\paragraph{No-cloning.}
No unitary clones an unknown state: $\braket{\phi}{\psi}=\braket{\phi}{\psi}^2$
forces overlaps $0$ or $1$. Basis states copy (classical bits);
superpositions do not. This underlies QKD security, the destruction of
the original in teleportation, and the impossibility of
copy-based error correction.

\paragraph{Teleportation and dense coding.}
Teleportation: $1\,\text{ebit}+2\,\text{cbits}\to1\,\text{qubit}$, with
Bob applying $\hat\sigma_z^{m_C}\hat\sigma_x^{m_A}$. Superdense coding,
its time-reverse: $1\,\text{ebit}+1\,\text{qubit}\to2\,\text{cbits}$ via
the four Bell states. No faster-than-light signalling: Bob's pre-message
state is $\tfrac12\id$.

\paragraph{Quantum key distribution.}
BB84 encodes bits in the conjugate $\hat\sigma_z$/$\hat\sigma_x$ bases;
E91 uses Bell pairs certified by CHSH. Intercept-resend injects a $25\%$
error rate; security rests on no-cloning and measurement disturbance,
not on computational hardness.

\paragraph{Clifford / magic dichotomy.}
Clifford circuits ($\{\mathrm{H},\mathrm{S},\mathrm{CNOT}\}$, Pauli prep
and measurement) are classically simulable in poly time
(Gottesman--Knill) via stabiliser tracking. Universality needs a
non-Clifford gate; $\mathrm{T}$ is supplied by distilled magic states
$\ket{\mathrm{T}}=\tfrac{1}{\sqrt2}(\ket{0}+e^{i\pi/4}\ket{1})$. By
Eastin--Knill no code has a universal transversal set, so magic states
dominate the cost of fault tolerance.

\end{chaptersummary}

% --------------------------------------------------------------
\section*{Exercises}
\addcontentsline{toc}{section}{Exercises}

\begin{exercise}[Hadamard basics]
Compute $\mathrm{H}\ket{0}$ and $\mathrm{H}\ket{1}$, and verify
$\mathrm{H}^2=\id$ and $\mathrm{H}\hat\sigma_x\mathrm{H}=\hat\sigma_z$.
\begin{solution}
$\mathrm{H}\ket{0}=\tfrac{1}{\sqrt2}(\ket{0}+\ket{1})=\ket{+}$ and
$\mathrm{H}\ket{1}=\tfrac{1}{\sqrt2}(\ket{0}-\ket{1})=\ket{-}$. Squaring,
$\mathrm{H}^2=\tfrac12\begin{psmallmatrix}1&1\\1&-1\end{psmallmatrix}
\begin{psmallmatrix}1&1\\1&-1\end{psmallmatrix}
=\tfrac12\begin{psmallmatrix}2&0\\0&2\end{psmallmatrix}=\id$. For the
conjugation, $\mathrm{H}\hat\sigma_x\mathrm{H}
=\tfrac12\begin{psmallmatrix}1&1\\1&-1\end{psmallmatrix}
\begin{psmallmatrix}0&1\\1&0\end{psmallmatrix}
\begin{psmallmatrix}1&1\\1&-1\end{psmallmatrix}
=\begin{psmallmatrix}1&0\\0&-1\end{psmallmatrix}=\hat\sigma_z$, the
first rule in \cref{eq:clifford-conjugation}.
\end{solution}
\end{exercise}

\begin{exercise}[All four Bell states from one circuit]
The circuit $\mathrm{CNOT}(\mathrm{H}\otimes\id)$ maps $\ket{00}$ to
$\ket{\Phi^+}$. What does it produce from the other three computational
basis inputs $\ket{01},\ket{10},\ket{11}$?
\begin{solution}
Apply $\mathrm{H}$ to the first qubit, then CNOT. From $\ket{01}$:
$\mathrm{H}\ket{0}\ket{1}=\ket{+}\ket{1}\to\tfrac{1}{\sqrt2}(\ket{01}+\ket{10})=\ket{\Psi^+}$.
From $\ket{10}$:
$\ket{-}\ket{0}\to\tfrac{1}{\sqrt2}(\ket{00}-\ket{11})=\ket{\Phi^-}$.
From $\ket{11}$:
$\ket{-}\ket{1}\to\tfrac{1}{\sqrt2}(\ket{01}-\ket{10})=\ket{\Psi^-}$.
The circuit is a unitary map from the computational basis onto the Bell
basis --- which is exactly the Bell measurement run backwards, as used
in superdense coding (\cref{stmt:superdense}).
\end{solution}
\end{exercise}

\begin{exercise}[CZ and CNOT]
Show $\mathrm{CZ}=(\id\otimes\mathrm{H})\,\mathrm{CNOT}\,(\id\otimes\mathrm{H})$,
and explain why CZ is symmetric between its two qubits while CNOT
appears not to be.
\begin{solution}
Conjugating the target of a CNOT by Hadamard turns the bit flip
$\hat\sigma_x$ into the phase flip $\hat\sigma_z$ (Exercise~1), so
$(\id\otimes\mathrm{H})\mathrm{CNOT}(\id\otimes\mathrm{H})
=\ket{0}\!\bra{0}\otimes\id+\ket{1}\!\bra{1}\otimes\hat\sigma_z=\mathrm{CZ}$.
CZ is diagonal, $\mathrm{diag}(1,1,1,-1)$, applying a phase only to
$\ket{11}$, which is manifestly symmetric under swapping the qubits.
CNOT looks asymmetric (a labelled control and target) but is
control/target-symmetric only up to Hadamards on both qubits; the phase
$-1$ on $\ket{11}$ is the gate's true symmetric core.
\end{solution}
\end{exercise}

\begin{exercise}[Cloning basis states but not superpositions]
A CNOT with the second qubit blank acts as $\ket{a}\ket{0}\to\ket{a}\ket{a}$
for $a\in\{0,1\}$. Show it fails to clone $\ket{+}$, and identify what it
produces instead.
\begin{solution}
By linearity,
$\mathrm{CNOT}\ket{+}\ket{0}
=\tfrac{1}{\sqrt2}(\mathrm{CNOT}\ket{0}\ket{0}+\mathrm{CNOT}\ket{1}\ket{0})
=\tfrac{1}{\sqrt2}(\ket{00}+\ket{11})=\ket{\Phi^+}$. A genuine clone
would give $\ket{+}\ket{+}=\tfrac12(\ket{00}+\ket{01}+\ket{10}+\ket{11})$,
a product state. Instead the output is the maximally entangled
$\ket{\Phi^+}$ --- the device copies the basis but entangles
superpositions, exactly as \cref{thm:no-cloning} requires.
\end{solution}
\end{exercise}

\begin{exercise}[Teleportation correction]
In the teleportation protocol Alice measures $(m_C,m_A)=(1,1)$. What is
Bob's state before correction, what does he apply, and why must he wait
for Alice's classical message?
\begin{solution}
From the derivation of \cref{stmt:teleportation}, outcome $(1,1)$ leaves
Bob in $\hat\sigma_x\hat\sigma_z\ket{\psi}=\alpha\ket{1}-\beta\ket{0}$.
He applies $\hat\sigma_z\hat\sigma_x$ (the inverse), since
$\hat\sigma_z\hat\sigma_x\,\hat\sigma_x\hat\sigma_z\ket{\psi}=\ket{\psi}$,
recovering the input. He must wait because, averaged over the four
equally likely outcomes, his reduced state is
$\tfrac14\sum_{m_C,m_A}\hat\sigma_x^{m_A}\hat\sigma_z^{m_C}\ket{\psi}\!\bra{\psi}\hat\sigma_z^{m_C}\hat\sigma_x^{m_A}=\tfrac12\id$,
which carries no information about $\ket{\psi}$ until he learns
$(m_C,m_A)$.
\end{solution}
\end{exercise}

\begin{exercise}[No signalling from teleportation]
Show explicitly that Bob's density matrix is independent of Alice's
input $\ket{\psi}$ until he receives the two classical bits, and explain
what this implies for faster-than-light communication.
\begin{solution}
Before the classical message, Bob's state is the equal mixture of the
four corrected branches,
$\op{\rho}_B=\tfrac12\id$, as computed in Exercise~5 --- the maximally
mixed state for \emph{any} $\ket{\psi}$. Since $\op{\rho}_B$ does not
depend on $\ket{\psi}$, no measurement Bob performs can reveal anything
about Alice's input. Information is transferred only when the two
classical bits arrive over the (sub-luminal) classical channel; the
entanglement alone transmits nothing, so there is no faster-than-light
signalling.
\end{solution}
\end{exercise}

\begin{exercise}[Decoding superdense coding]
In superdense coding Alice wants to send the bits $(1,0)$. Which Bell
state does she create, and what sequence of operations lets Bob read off
the bits?
\begin{solution}
For $(b_1,b_2)=(1,0)$ Alice applies
$\hat\sigma_z^{1}\hat\sigma_x^{0}=\hat\sigma_z$ to her qubit, turning
$\ket{\Phi^+}$ into $\ket{\Phi^-}=\tfrac{1}{\sqrt2}(\ket{00}-\ket{11})$.
After receiving Alice's qubit, Bob applies $\mathrm{CNOT}_{A\to B}$ then
$\mathrm{H}$ on $A$ and measures both qubits: this Bell measurement maps
the four Bell states back to the computational basis, returning
$\ket{10}$ and hence $(b_1,b_2)=(1,0)$ with certainty, since the Bell
states are orthogonal.
\end{solution}
\end{exercise}

\begin{exercise}[The 25\% error rate in BB84]
Derive the bit-error rate that an intercept-resend eavesdropper injects
into the sifted key of BB84, and state how Alice and Bob detect her.
\begin{solution}
On a sifted round Alice and Bob used the same basis. Eve, ignorant of
it, guesses the basis correctly half the time (no disturbance) and
wrongly half the time. On a wrong-basis round she measures an eigenstate
of the conjugate observable and resends it; Bob's measurement in the
correct basis then yields the wrong bit with probability $\tfrac12$
(e.g.\ $|\!\braket{-}{0}\!|^2=\tfrac12$). The induced error rate is
$\tfrac12\times\tfrac12=25\%$. Alice and Bob publicly compare a random
subset of sifted bits; an error rate near $25\%$ (well above the
$\approx11\%$ security threshold and any honest channel noise) reveals
the eavesdropper, and the key is discarded.
\end{solution}
\end{exercise}

\begin{exercise}[The Clifford status of $\mathrm{S}$ and $\mathrm{T}$]
Decide whether $\mathrm{S}$ and $\mathrm{T}$ are Clifford by conjugating
$\hat\sigma_x$, and state the consequence for universality.
\begin{solution}
$\mathrm{S}\hat\sigma_x\mathrm{S}^\dagger
=\begin{psmallmatrix}1&0\\0&i\end{psmallmatrix}
\begin{psmallmatrix}0&1\\1&0\end{psmallmatrix}
\begin{psmallmatrix}1&0\\0&-i\end{psmallmatrix}
=\begin{psmallmatrix}0&-i\\i&0\end{psmallmatrix}=\hat\sigma_y$, a Pauli
--- so $\mathrm{S}$ is Clifford. But
$\mathrm{T}\hat\sigma_x\mathrm{T}^\dagger
=\begin{psmallmatrix}0&e^{-i\pi/4}\\e^{i\pi/4}&0\end{psmallmatrix}
=\tfrac{1}{\sqrt2}(\hat\sigma_x+\hat\sigma_y)$, which is \emph{not} a
Pauli --- so $\mathrm{T}$ is non-Clifford. By
\cref{stmt:gottesman-knill} the Clifford set is classically simulable;
adding $\mathrm{T}$ is what lifts the gate set to universality.
\end{solution}
\end{exercise}

\begin{exercise}[Why magic states are unavoidable]
Using the Eastin--Knill theorem, explain why a fault-tolerant computer
based on a stabiliser code cannot implement \emph{every} gate
transversally, and how magic states resolve the impasse.
\begin{solution}
A transversal gate acts independently on corresponding physical qubits
of code blocks, so a single physical fault cannot spread within a block
--- the ideal fault-tolerant implementation. Eastin--Knill states that
no nontrivial code admits a transversal set that is also universal: the
transversal gates always form a finite (in practice, Clifford-like)
group, which \cref{stmt:gottesman-knill} shows is non-universal. At
least one gate, conventionally $\mathrm{T}$, must therefore be
implemented some other way. Magic-state distillation
(\cref{stmt:magic-distillation}) supplies it: high-fidelity
$\ket{\mathrm{T}}$ states are prepared offline by Clifford-only circuits
and injected through the $\mathrm{T}$ gadget, confining the
non-transversal cost to a magic-state factory while keeping every
in-line operation Clifford and fault-tolerant.
\end{solution}
\end{exercise}
